\section{Adaptive blocking flows for cocycle-span minimization}
\label{sec:cocycle-blocking}

The span network has one supply unit per tetrahedron. The previous solver
runs Dijkstra once per unit, even when many units can travel at the same
reduced cost. The maintained optimizer now routes batches through
zero-reduced-cost residual arcs, using shortest-path searches when the
current zero-cost graph is exhausted or batching has been unproductive.
Its final primal--dual certificate and independent verifier are unchanged.
This improves a polynomial local primitive; it does not complete the
normal-surface search or establish a general quasipolynomial recognizer.

\subsection{Routing a batch while retaining dual feasibility}

Recall the $O(T)$-node, $O(T)$-arc transshipment network of
Section~\ref{sec:cocycle-seeds}. Each lower node receives at most one
unit from the source, each upper node sends at most one unit to the sink,
and every internal forward arc has capacity $T+1$. An acyclic shortest-path
initialization supplies integral potentials $\pi$ with nonnegative
reduced costs
\[
 r(u,v)=c(u,v)+\pi(u)-\pi(v)
\]
on every usable residual arc.

An augmentation entirely along arcs with $r=0$ needs no change of
potential. Every newly opened reverse arc also has reduced cost zero;
every other usable arc retains its previous nonnegative reduced cost.
Feasibility of the potentials is therefore preserved after any number of
such augmentations. A path has unit bottleneck because its first source
arc has capacity one. The numerical size of a height does not determine
the number of flow units.

The producer first assigns breadth-first levels in the usable zero-cost
graph. Once the sink level is known, it finishes the preceding layer and
does not expand deeper vertices. It then traverses only arcs advancing
one level. A current-arc pointer at each node skips exhausted or unusable
choices. Traversal uses explicit node and edge stacks, so a long residual
path does not depend on the interpreter's recursion limit.

When a path reaches the sink, the producer sends one unit and restarts
at the source. A failed branch advances its parent's pointer. The
resulting block leaves no source-to-sink path in the current level graph,
unless all required flow has already been sent. Newly opened reverse arcs
go to a lower level and cannot revive an earlier discarded choice in the
same block. Further zero-cost paths, if any, are found by another
breadth-first level construction. This argument allows zero-cost cycles,
parallel arcs and disconnected abstract incidence components.

If no zero-cost path remains, reduced-cost Dijkstra finds the next path.
Let $\Delta$ be the sink distance. The existing update
\[
 \pi(v)\gets\pi(v)+\min(d(v),\Delta)
\]
uses increment $\Delta$ for unreachable or unsettled nodes. It preserves
nonnegative residual costs, including on arcs entering the reached set
from outside, and makes the chosen augmenting path tight. The producer
sends one unit along that path, then may try batching again.

After $T$ units have been sent, every original internal forward arc is
still residual: its capacity exceeds the total flow. Every used arc also
has its reverse residual, so dual feasibility forces equality on it.
Extracting vertex potentials and a corner matching therefore gives the
same primal--dual optimality argument as before. The existing independent
checker reconstructs coordinates, verifies matching balance and requires
equality of the primal span and dual reward. It neither imports the
blocking-flow helper nor trusts its counters.

\subsection{A bounded handoff for small batches}

Batch construction itself costs a graph scan. When marginal costs are
distinct, repeating that scan before every Dijkstra search adds work
without routing more flow. The scheduler therefore counts an attempted
block as poor when it sends fewer than two units. Two consecutive poor
attempts disable batching; a block of at least two units resets that count.
The solver then uses ordinary shortest-path augmentations.

A later Dijkstra augmentation with sink distance zero is fresh evidence
that units are again moving at the same reduced-cost level. It reenables
batching with a reset poor-attempt counter. This restart affects only
which valid augmenting paths are tried. Both modes maintain the same
residual feasibility and final certificate obligations. Every iteration
of the outer loop sends at least one unit, including an iteration whose
zero-cost search fails, because that failure is followed by Dijkstra.

The separated-cost control has $T$ independent vertex labels, one for
each row, with heights $(0,0,0,i)$ in row $i$. All local spans are fixed,
and successive marginal costs are distinct. The implementation attempts
only two zero-cost searches, sends one unit in its first block and then
finishes with $T-1$ Dijkstra searches. The unbounded prototype instead
performed a failed zero-cost scan at almost every remaining unit; this
control motivated the handoff before integration.

\subsection{Worst-case accounting and an improved equal-cost family}

Write $N,E=O(T)$ for the network size. One breadth-first level construction
uses $O(N+E)$ graph work. Within its block, current-arc pointers advance
at most once per arc. A successful path can traverse at most $N-1$ arcs;
if a block sends $q$ units, a conservative cost is
$O(N+E+qN)$. Across the whole solve there are at most $T$ attempted
blocks and at most $T$ routed units. Thus all blocking work is
$O(T^2)$ arithmetic and incidence work. At most $T$ Dijkstra searches
remain, preserving the conservative bound
\[
 O(T^2\log(T+1))
\]
for arithmetic operations. Arrays, queues, current-arc pointers and
explicit stacks take $O(T)$ records. The handoff changes neither bound.

Zero-cost augmentations do not change potentials. The previous bound
$O(B+\log(T+1))$ on arithmetic bit length for $B$-bit heights therefore
still applies to both modes. A conservative addition/comparison bit bound
remains
\[
 O\!\left(T^2\log(T+1)(B+\log(T+1))\right),
\]
apart from label indexing and input/output. The original certificate
size and separate cohomology-preprocessing bounds are unchanged.

There is a direct improvement on an explicit family. Give every row
vertices $(0,1,2,3)$ and heights $(0,0,0,0)$. The zero-cost level graph
has source, lower, vertex, upper and sink layers. One blocking phase sends
all $T$ units along paths of length four, with $O(T)$ scans and no
Dijkstra search. The old solver performs $T$ searches, each scanning the
$T$ source arcs, so its residual search work is $\Omega(T^2)$ even on
this zero-height family. The change is from quadratic to linear residual
search work on that family, not a linear bound for arbitrary span inputs
or complete unknot recognition. Complete-call measurements below include
certificate creation and independent verification as well.

\subsection{Arithmetic, topology and source audits}

The 27 focused blocking, span, seed and tree tests pass in 2.593 seconds.
The full suite passes all 1,192 tests in 212.108 seconds. New tests cover
linear scan growth on the equal-cost family, bounded handoff on distinct
marginal costs, restarting after zero-distance evidence, a 2,500-node
zero-cost path with cycles, cooperative caps and cancellation. Multiplying
a nonzero instance by $2^{20000}$ preserves all search counters and scales
the objective exactly. Its hexadecimal JSON certificate replays with both
flow producer entry points disabled.

Five hundred random arithmetic problems agree with the old optimizer on
both optimum value and the entire coordinate vector. In 495 cases the
entire certificate is also identical; five use another valid witness.
This observed coordinate agreement is not assumed in the
correctness proof. Each old and new certificate passes the unchanged
independent checker. Aggregate Dijkstra searches fall from 7,928 to 7,398,
but total guard work rises from 3,333,719 to 3,765,178. Small arbitrary
inputs therefore provide an adverse control, not evidence of a uniform
speedup. The complete-call benchmark includes a separate batch of such
problems.

The initial audit completed its checks but failed while serializing the
20,000-bit objective with ordinary decimal JSON. The driver now uses the
existing exact hexadecimal transport; the interpreter's conversion limit
is unchanged. The final audit is repeated after that fix and takes
51.768 seconds. Its elapsed time mixes correctness checks and is not a
paired performance comparison. The initial failure log is retained.

All 84 source inputs are also optimized directly with a two-million-guard
allowance. Both optimizers finish on 82 cases, preserving every objective
and coordinate vector. Six complete source span witnesses are byte-identical;
the others can use different valid matchings. The new optimizer also
finishes the two previously capped standalone queries:
\begin{itemize}
\item Gordian, with 2,820 tetrahedra, uses six Dijkstra searches and routes
      2,814 units in blocks. It reduces 3,513 pieces to 2,315 in 1,929,649
      guards. An independent source-bound inspection and Regina agree
      that its surface is connected, orientable and has $\chi=-65$,
      with no essential disc.
\item The 32-crossing circle uses 18 Dijkstra searches for 640 units,
      reducing 8,385 pieces to 4,235 in 1,019,747 guards. Independent
      inspection and Regina identify a disc. Its raw surface already
      supplied a native proof before this change.
\end{itemize}
These two larger topology checks supplement 74 fresh optimized-surface
Regina controls on sources of at most twelve crossings. All 76 comparisons
agree on connectedness, orientability, Euler characteristic and disc count.

Complete source-stage queries preserve the four-tree policy's 23 positives,
60 completed restricted misses and one cap, and the 24-tree policy's 24
positives, 59 misses and one cap. Gordian still caps because source
construction, cohomology, tree trials and replay share its two-million-guard
allowance with optimization. Its newly completed standalone query is not
a newly completed full stage. The selected optimized positive and its
duplicate \code{random-54} now have a different valid corner matching,
giving one distinct changed positive proof. Other positive hashes are
preserved, and every returned source proof replays independently.

The default-stage guards for \code{survivor-00}, \code{survivor-01} and
\code{survivor-02} fall from 1,310,623, 1,124,340 and 1,234,216 to
251,349, 208,359 and 172,493 respectively. These remain restricted misses.
The optimized five-crossing positive falls from 144,578 to 54,318 guards,
and the optional late-tree positive falls from 180,578 to 72,527. Guard
counts are cooperative work indicators, not bit-operation counts.

On nineteen ordinary corpus inputs, old and new enabled-stage policies
return the same expected verdicts without reaching this native stage.
The optimization improvements therefore do not establish an ordinary
default-portfolio speedup or new source coverage.

\subsection{Complete optimization and recognition measurements}

The final measurements run serially, without tests, audits, builds or
commits overlapping them. Each case has one warm-up round followed by
five measured rounds, each randomly ordering old, old A/A, new and new
A/A endpoints. All 260 measured endpoints and 52 warm-ups complete.
Tables give median primary-arm times and medians of paired ratios;
the latter need not equal ratios of the median times. Each implementation's
certificate hashes remain stable across repetitions, and all comparable
coordinate hashes agree between implementations. Valid matching choices
can make complete certificate hashes differ.

The first experiment includes input construction, span optimization,
an additional independent certificate check and hashing. For the two
source-derived cases, every timed endpoint also constructs the canonical
exterior and solves its primitive cohomology. The binary case scales its
four-crossing cocycle heights by $2^{4096}$ before making the arithmetic
span query; it is not being used as a primitive source-disc certificate.
The random control aggregates 24 complete, independently verified
24-row problems per endpoint to reduce timer noise on small queries.
The other optimizer endpoints contain one query each. This mode completes
140 measured endpoints and 28 warm-ups in 46.753 seconds.

\begin{center}
\small\begin{tabular}{lrrrrr}
\toprule
Input & Old (ms) & New (ms) & Old/new & Old A/A & New A/A \\
\midrule
Equal cost, 32 rows & 5.594 & 1.964 & 2.849 & 1.006 & 0.999 \\
Equal cost, 128 rows & 75.118 & 8.017 & 9.064 & 1.012 & 1.001 \\
Equal cost, 512 rows & 1164.347 & 33.169 & 35.563 & 1.007 & 1.025 \\
Distinct costs, 128 rows & 74.723 & 74.804 & 0.985 & 1.004 & 0.999 \\
24 random 24-row queries & 111.174 & 130.429 & 0.858 & 1.001 & 1.069 \\
Circle, 32 crossings & 1683.749 & 429.510 & 3.914 & 1.000 & 1.015 \\
Circle, 4 crossings, binary & 49.438 & 27.429 & 1.811 & 1.008 & 0.988 \\
\bottomrule
\end{tabular}

\end{center}

On the equal-cost family, complete endpoints improve from $2.85\times$
at 32 rows to $35.56\times$ at 512 rows, consistent with the proved
residual-search improvement. The largest median falls from 1.164 seconds
to 0.033 seconds. The distinct-cost control is close to parity, with a
paired ratio of 0.985 after the bounded handoff. The random control is
about 17 percent slower by its paired ratio; its new-arm A/A ratio of
1.069 also shows visible variation. This is consistent with the adverse
guard-work result above, and prevents a claim of uniformly faster span
optimization.

The two source-derived optimization endpoints improve $3.91\times$ and
$1.81\times$. The 32-crossing case includes source construction and
cohomology, falling from 1.684 to 0.430 seconds. This is a timing of its
optimization query: the recognizer normally succeeds on its raw surface
and does not run that optimization.

The second experiment starts from the input PD and times the complete
recognizer, including its earlier setup, native discovery, exact fallback
when required, an additional independent replay of any returned native
proof and certificate/coordinate hashing. Earlier shortcuts are disabled
to make the native stage reachable. Both arms use no local seed cap and
retain the recognizer's twelve-second global deadline. The late-tree
case requests 24 trials; the others use four. This mode completes 120
measured calls and 24 warm-ups in 11.591 seconds.

\begin{center}
\small\begin{tabular}{lrrrrr}
\toprule
Input & Old (ms) & New (ms) & Old/new & Old A/A & New A/A \\
\midrule
Raw circle & 63.035 & 61.919 & 1.034 & 1.032 & 0.995 \\
Optimized positive & 128.244 & 89.515 & 1.459 & 0.956 & 0.974 \\
Early tree positive & 37.066 & 38.425 & 0.999 & 1.000 & 1.005 \\
Late tree positive & 178.053 & 139.721 & 1.280 & 0.964 & 1.049 \\
Genus-one miss & 57.538 & 41.679 & 1.378 & 0.997 & 0.986 \\
Trefoil & 55.427 & 39.291 & 1.404 & 0.976 & 1.000 \\
\bottomrule
\end{tabular}

\end{center}

The four cases reaching span optimization improve $1.28$--$1.46\times$.
They preserve the optimized and late-tree positive proofs' conclusions,
and preserve the genus-one unknot miss and trefoil's exact fallback
verdicts. The early tree positive has a paired ratio of 0.999 and avoids
flow in both arms. The raw-circle control's 1.034 ratio is comparable to
its old-arm A/A ratio of 1.032; no algorithmic gain is claimed for that
unchanged branch. These are selected native-stage timings, not evidence
of improved ordinary portfolio performance or a general recognition bound.


Reproduction uses \path{../fast/normal_orbit_research/blocking.py}, with
\code{audit}, \code{optimize} and \code{recognize} modes. The baseline
is \code{4fa723b79}; all 483 runtime, driver, fixture and prior-audit
source pins agree. The two large topology checks use
\path{data/blocking_capped_sources.py}, run from \code{fast} with
\code{PYTHONPATH=.} and optional Regina. The incoming archive inventory
and the checked \code{openai/math} head \code{fd4aeeb2e} remain unchanged.

